\section{Literature Review}

This section provides an integrated review of the literature related to judicial mandates and public procurement, with a particular focus on the health sector. The review is organized into several subsections corresponding to key dimensions of our study: the impact of judicial mandates, the institutional context, theoretical frameworks, empirical evidence, and policy implications. We also highlight the research gaps that our paper addresses.

Judicial mandates have been increasingly scrutinized in the literature as a critical factor influencing public procurement processes. Early works highlight the role of courts in protecting individual rights by ensuring access to essential services, particularly in health care. However, subsequent studies have revealed that such interventions, while well-intentioned, may disrupt standard procurement mechanisms. Researchers have noted that court orders often force governments into urgent purchasing scenarios, reducing the time available for proper planning and competitive bidding.

Several studies, including \citet{wang2015right} and \citet{cnj2019judicializacao}, document how judicial interventions lead to deviations from normal procurement practices. These interventions typically result in acquisitions being made under time constraints and with fewer suppliers participating, which in turn can drive up prices. The literature emphasizes that while judicial mandates aim to guarantee timely access to critical health services, they inadvertently impose additional costs on public budgets.

Furthermore, recent research has begun to explore the so-called “under the gun” effect, whereby the threat of punitive sanctions further elevates procurement costs. However, existing studies often treat judicial interventions as a binary phenomenon, without dissecting the nuances between different types of urgent procurements—namely, those that are purely administrative versus those that are judicially enforced. This is one of the key gaps our paper addresses.

In summary, while the literature provides ample evidence that judicial mandates increase procurement costs, there remains a lack of formal models that explicitly quantify the welfare losses associated with these disruptions. Our paper fills this gap by developing a theoretical framework that distinguishes between ordinary, urgent, and litigated procurement regimes, thus offering a more granular understanding of the cost implications.

The institutional context in which public procurement occurs is fundamental to understanding its performance. In Brazil, the Unified Health System (SUS) represents a landmark institutional framework intended to guarantee universal health coverage. However, as noted by \citet{castro2019} and \citet{soares2019}, the rigid interpretation of constitutional rights in this system often forces public agencies to operate under extreme urgency. This urgency, in turn, distorts procurement processes and strains public budgets.

Institutional challenges are not limited to the health sector. Broader literature in public administration, such as \citet{brewer2010impact} and \citet{basheka2011economic}, discusses the pervasive issues of bureaucratic inefficiencies and red tape that exacerbate the difficulties in executing well-planned public procurement. These studies highlight that administrative constraints—such as inadequate planning and lack of inter-agency coordination—compound the adverse effects of judicial mandates.

Moreover, while several works examine the operational challenges within public agencies, there is still limited discussion on how these institutional shortcomings interact with judicial pressures to alter market outcomes in procurement. Our study contributes to this debate by linking institutional rigidity with the increased costs observed in urgent procurement processes.

Finally, by integrating insights from both public administration and judicial studies, our paper addresses the gap in understanding how the interplay between legal mandates and bureaucratic inefficiencies can lead to significant welfare losses. This integrated perspective is crucial for designing effective policy interventions that reconcile the need for timely health interventions with fiscal responsibility.

The development of theoretical frameworks to analyze public procurement under judicial pressure has gained momentum over the past decades. Foundational work in auction theory has provided robust tools to study competitive bidding processes, particularly in settings where external shocks alter traditional market dynamics. Seminal papers such as \citet{bandiera2009active} have laid the groundwork by distinguishing between active and passive waste in government spending, offering a lens through which the effects of urgency on procurement costs can be understood.

Building on these foundations, our theoretical model incorporates elements from both auction theory and welfare economics. The model distinguishes between three procurement regimes—ordinary, urgent, and litigated urgent—each characterized by distinct levels of supplier participation and competitive intensity. Although prior studies (e.g., \citet{coviello2017court}) have examined judicial impacts on procurement, they typically do not formalize the role of judicial penalties as a stochastic element within an auction framework. This omission represents a significant gap in the literature that our paper seeks to fill.

Moreover, literature from the field of strategic management provides valuable insights into how external pressures influence organizational behavior. Works by \citet{barney1991firm}, \citet{eisenhardt1989making}, and \citet{grant1996toward} illustrate that competitive dynamics and decision-making processes are sensitive to environmental shocks. These studies support our argument that judicial mandates not only affect prices but also alter the strategic behavior of both suppliers and public buyers.

Lastly, by integrating theoretical insights from auction theory and strategic management, our model offers a novel perspective on the “under the gun” effect—a relatively underexplored aspect in the current literature. This contribution is particularly relevant for understanding how judicial sanctions can exacerbate procurement inefficiencies beyond the mere reduction in competition.

Empirical studies on public procurement have increasingly utilized detailed administrative datasets to uncover the inefficiencies inherent in public purchasing processes. A number of studies, including those by \citet{mironov2016corruption} and \citet{lewisfaupel2016can}, demonstrate that factors such as reduced competition and strategic supplier behavior are crucial determinants of procurement costs. Despite this progress, there remains a gap in the empirical literature regarding the differential impacts of judicial versus administrative urgent procurements.

The use of granular data allows researchers to distinguish between different procurement regimes and to measure the specific cost increments associated with urgent and litigated procurement processes. Empirical evidence indicates that when procurement is conducted under urgency—especially in cases where judicial mandates are involved—the resulting prices are significantly higher. However, many studies aggregate all urgent procurements into a single category, thereby masking the distinct effects of judicial penalties. Our paper addresses this gap by disaggregating urgent procurements into administrative and litigated types, allowing for a more refined analysis of cost drivers.

Furthermore, empirical research has shed light on the broader implications of inefficient procurement practices. Studies such as those by \citet{cullen2016performance} and \citet{moore2012cartels} reveal that procurement inefficiencies not only inflate costs but also lead to resource misallocation and reduced overall welfare. Nevertheless, the quantification of welfare losses specifically attributable to judicial mandates remains underexplored. Our work contributes to this debate by providing both theoretical and empirical estimates of these welfare losses.

Finally, cross-country comparisons in the empirical literature suggest that the negative effects of urgent procurement are not unique to Brazil. Comparative studies, such as those examining procurement in India and Indonesia \citep{lewisfaupel2016can}, highlight similar patterns of cost escalation under urgency. This broader perspective reinforces the importance of our study, which seeks to generalize the findings and contribute to a more comprehensive understanding of the phenomenon across different institutional settings.

The convergence of theoretical and empirical evidence on judicial mandates and procurement inefficiencies has important policy implications. Researchers such as \citet{perry1988public} and \citet{otoolle2003bureaucracy} have long argued for improved coordination between public agencies and oversight bodies to enhance procurement efficiency. Nonetheless, there remains a gap in understanding how specific judicial interventions translate into measurable welfare losses and what policy levers can be used to mitigate these adverse effects.

One major policy implication is the need for enhanced coordination between the judiciary and the executive branches. By reducing the frequency and severity of judicial interventions, it may be possible to alleviate the “under the gun” effect and restore competitive equilibrium in procurement auctions. Our study contributes to this debate by quantifying the cost differentials between administrative and litigated procurements, thereby providing concrete targets for policy reform.

Moreover, revising penalty mechanisms to avoid excessively punitive measures could help reduce the negative impact on procurement costs. Research in public administration, such as the work by \citet{williams2018management} and \citet{rasul2018management}, suggests that more balanced penalty structures may incentivize better planning and execution without compromising accountability. Despite these insights, the optimal design of such penalties remains an open question—a gap that our analysis begins to address.

Finally, future research should extend this line of inquiry by examining the impact of judicial mandates across different institutional settings and sectors. Comparative analyses can provide deeper insights into the structural factors that mediate the relationship between judicial intervention and procurement inefficiency. Our work lays the groundwork for such studies by establishing a theoretical framework and providing initial empirical estimates of welfare losses, thereby contributing to a more nuanced understanding of the issue and informing future policy directions.
